\documentclass[11pt,a4paper]{article}
\usepackage[T1]{fontenc}
\usepackage{lmodern}
\usepackage[margin=26mm,headheight=14pt]{geometry}
\usepackage{amsmath,amssymb,amsthm,mathtools}
\usepackage{microtype,booktabs,longtable,array}
\usepackage{xcolor,xurl}
\definecolor{linkink}{RGB}{25,55,85}
\usepackage[colorlinks=true,allcolors=linkink,pdfusetitle]{hyperref}
\usepackage{fancyhdr}
\pagestyle{fancy}\fancyhf{}
\fancyhead[L]{\small Quadratic polylogarithm boundaries}
\fancyhead[R]{\small ProveIt research contribution}
\fancyfoot[C]{\thepage}
\newtheorem{theorem}{Theorem}[section]
\newtheorem{proposition}[theorem]{Proposition}
\newtheorem{lemma}[theorem]{Lemma}
\newtheorem{corollary}[theorem]{Corollary}
\theoremstyle{definition}\newtheorem{definition}[theorem]{Definition}
\theoremstyle{remark}\newtheorem{remark}[theorem]{Remark}
\newcommand{\Li}{\operatorname{Li}}
\newcommand{\ee}{\mathrm e}
\newcommand{\ii}{\mathrm i}
\newcommand{\dd}{\,\mathrm d}
\newcommand{\C}{\mathbb C}
\newcommand{\R}{\mathbb R}
\newcommand{\N}{\mathbb N}
\newcommand{\Q}{\mathbb Q}
\newcommand{\Z}{\mathbb Z}
\newcommand{\Uc}{U_{\mathrm c}}
\newcommand{\vb}{v_{\mathrm b}}
\newcommand{\coef}[2]{[#1]#2}
\newcommand{\Dop}{\mathcal D}
\numberwithin{equation}{section}
\allowdisplaybreaks[2]
\title{Critical Boundaries of Quadratic\\Polylogarithm Feedback\\[5pt]
\large Exact boundary transfer, convergent Hahn inverses,\\
Gamma cancellations, and the marginal Lambert law}
\author{Research contribution prepared for Vladimir Reshetnikov\\
\small with OpenAI assistance}
\date{12 October 2026}
\hypersetup{pdftitle={Critical Boundaries of Quadratic Polylogarithm Feedback},pdfauthor={Research contribution prepared for Vladimir Reshetnikov with OpenAI assistance},pdfsubject={Polylogarithms, boundary transfer, convergent Hahn inverses and Lambert logarithmic asymptotics}}
\begin{document}
\maketitle
\begin{abstract}
We resolve the fixed-parameter boundary problem explicitly left open in the
incoming article \emph{Resonant Polylogarithm Transseries}. For the analytic
germ defined by
\(uv-v^2/2=\Li_s(r\ee^v)\), we prove its exact convergence radius and unique
dominant singularity at and beyond
\(u=\sqrt{2\zeta(s)+\zeta(s-1)^2}\), for every real $s>2$.
Beyond criticality we prove the conjectured $R^{-n}n^{-s}$ coefficient law,
including its constant and every positive integer order. At criticality we
construct convergent two-generator inverse charts and obtain the fractional,
logarithmic, and finite-curvature regimes.

The logarithmic endpoint $s=3$ has an explicitly normalized Lambert core.
A finite coefficient formula, followed by reciprocal-Gamma differentiation,
gives every inverse-logarithmic correction to its coefficient asymptotics;
the first two corrections are displayed in closed form. A polylogarithmic
transfer dictionary eliminates spurious descendants of logarithmic-coordinate
powers. At $s=5/2$ this proves an exact missing $n^{-2}$ term. At $s=4$ it
instead produces a nonzero $n^{-2}$ correction and proves that the leading
square root is not a genuine Puiseux germ. We classify finite ramification
at all noninteger critical orders and extend the boundary constants and
marginal normalization to finite changes of the coefficient prefix.

All theorems have conventional proofs. Exact finite checks, positive
coefficient recurrences, numerical diagnostics, an integration crosswalk,
and further research questions accompany the article. Classical inversion,
singularity analysis, and simply generated tree mechanisms are credited;
worldwide priority for the model-specific refinements is not asserted.
\end{abstract}
\medskip
\noindent\textbf{Status.} Unrefereed research manuscript, not a proof-assistant
formalization. Numerical evidence is not an interval certificate. The
repository observation and limits of the source audit are recorded in
Section~\ref{qb:sec:audit}; no remote repository file was changed.
\clearpage
\tableofcontents
\clearpage

\section{The research question and what is added}\label{qb:sec:intro}

The starting point is the quadratic feedback equation
\begin{equation}\label{qb:eq:model}
 uv(r)-\frac{v(r)^2}{2}=\Li_s\!\left(r\ee^{v(r)}\right),
 \qquad v(0)=0,\qquad v(r)=\sum_{n\ge1}a_n r^n.
\end{equation}
The parameters $s$ and $u$ are real. Our main boundary theorems assume
$s>2$ and $u\ge\Uc(s)$, where
\begin{equation}\label{qb:eq:critical}
 A=\zeta(s),\qquad B=\zeta(s-1),\qquad
 \Uc(s)=\sqrt{2A+B^2}.
\end{equation}
The variable $r$ is complex near zero. Substituting $r=\ee^{-x}$ converts
its Taylor series into a convergent exponential-sector expansion. Near its
first singularity a different inversion problem appears, with fractional
powers or logarithms of the distance to the boundary. These two levels are
why the model belongs simultaneously to polylogarithm analysis and to the
repository's transseries inversion program.

Equation~\eqref{qb:eq:model} is a \emph{quadratic feedback model}. Except at
the dilogarithm specialization discussed in the predecessor, it is not the
exact-core inversion equation for $-\Li_s(-\ee^x)$. The distinction is
important when comparing this article with general integer-order Fermi
inversion.

\subsection{A precise pending question}

The incoming article \cite{prior} proves the interior-fold regime. Its
Section~12, Question~1 explicitly leaves the case $s>2$, $u\ge\Uc(s)$ open.
For $u>\Uc(s)$ it proposes
\begin{equation}\label{qb:eq:conjecture}
 a_n\sim \frac{w^{s-1}}{(w-B)^s}R^{-n}n^{-s},\qquad
 w=\sqrt{u^2-2A},\quad \vb=u-w,\quad R=\ee^{-\vb}.
\end{equation}
It correctly notes that formal transport of a singular term does not prove
this assertion: one must establish complex continuation, treat resonant
integer orders, and exclude other dominant singularities. We supply these
steps and prove \eqref{qb:eq:conjecture}. We also determine the critical
cases anticipated there.

\begin{theorem}[Fixed-parameter boundary classification]\label{qb:thm:summary}
Let $s>2$ and $u\ge\Uc(s)$ in \eqref{qb:eq:model}. With $w,\vb,R$ as in
\eqref{qb:eq:conjecture}, the coefficients are strictly positive, the exact
radius is $R$, and $r=R$ is the only singularity on $|r|=R$. The series is
absolutely convergent on that circle and $v(R)=\vb$.

For $u>\Uc(s)$, \eqref{qb:eq:conjecture} holds. For $u=\Uc(s)$, so that
$w=B$, the following laws hold as $n\to\infty$:
\begin{align}
 a_n R^n&\sim
 \frac{(B/\Gamma(1-s))^{1/(s-1)}}{-\Gamma(-1/(s-1))}
 n^{-1-1/(s-1)}, &&2<s<3,\label{qb:eq:summary-stable}\\
 a_n R^n&\sim \sqrt{\frac{B}{\pi}}\,
 n^{-3/2}(\log n)^{-1/2}, &&s=3,\label{qb:eq:summary-log}\\
 a_n R^n&\sim
 \sqrt{\frac{B}{2\pi(1+\zeta(s-2))}}\,n^{-3/2}, &&s>3.
 \label{qb:eq:summary-finite}
\end{align}
These are fixed-parameter statements, not uniform transition laws as $s$
or $u$ varies.
\end{theorem}

The proof is assembled in Sections~\ref{qb:sec:global}--\ref{qb:sec:finite}.
The constants in \eqref{qb:eq:summary-stable} are positive: for
$1<\alpha=s-1<2$, $\Gamma(-\alpha)>0$ and
$\Gamma(-1/\alpha)<0$.

\subsection{Antecedents and proposed contribution}

Polylogarithm continuation and its Gamma--zeta expansion are classical
\cite{dlmfpoly}. So are Lagrange inversion, the positive implicit-function
schema, Hankel transfer, and logarithmic singularity analysis
\cite{fo,fs}. The substitution used below places the problem in the simply
generated tree/compound-Poisson setting surveyed by Janson \cite{janson}.
The repository already has a linear-core critical Hahn theory, with
ramification and arithmetic cancellation results, and related marginal
and confluent packages \cite{repocritical,repotrans}. Those mechanisms are
not discoveries of this article.

The proposed additions are narrower and explicit: completion of the
predecessor's quadratic boundary theorem; its all-order two-variable
charts on both sides of $s=3$; the exact marginal Lambert normalization and
finite Gamma-derivative coefficient compiler; the $s=5/2$ and $s=4$
coefficient identities; and a finite-prefix dictionary that tracks the
quadratic correction instead of importing linear-core constants unchanged.
The explicit square-root lift below explains both the connection with
classical theory and the extra constants it introduces.

\section{The positive quadratic lift and exact coefficients}\label{qb:sec:lift}

Define, with the square root positive at $t=0$,
\begin{equation}\label{qb:eq:lift}
 h(t)=u-\sqrt{u^2-2\Li_s(t)},\qquad t=r\ee^{h(t)},\qquad v(r)=h(t(r)).
\end{equation}
Since $h(0)=0$ and $u>0$, the
analytic implicit function theorem gives a unique germ $t(r)$, and hence
$v(r)$, near zero.

\begin{proposition}[Three exact coefficient descriptions]\label{qb:prop:coeff}
For any real $s$ and $u>0$, let $h(t)=\sum_{j\ge1}h_jt^j$ be the lift germ.
Then
\begin{align}
 u h_j&=j^{-s}+\frac12\sum_{k=1}^{j-1}h_kh_{j-k},\label{qb:eq:hrec}\\
 a_n&=\frac1n[t^n]\exp(nh(t)),\label{qb:eq:lagrange}\\
 a_n&=\sum_{k=1}^n\frac{C^{(s)}_{n,k}}{k}
 \sum_{j=0}^{k-1}\binom{k+j-1}{j}
 \frac{n^{k-1-j}}{(k-1-j)!\,2^j u^{k+j}},\label{qb:eq:finite-coeff}\\
 C^{(s)}_{n,k}&=[z^n]\Li_s(z)^k
 =\sum_{\substack{n_1+\cdots+n_k=n\\n_j\ge1}}
 (n_1\cdots n_k)^{-s}.\notag
\end{align}
In particular $h_j>0$ and $a_n>0$. Formula~\eqref{qb:eq:finite-coeff} is the
predecessor's positive finite formula, restated for comparison rather than
claimed anew.
\end{proposition}
\begin{proof}
The identity $uh-h^2/2=\Li_s(t)$ gives \eqref{qb:eq:hrec}, and induction
proves positivity. Lagrange inversion in $t=r\ee^{h(t)}$ gives
\[
 [r^n]h(t(r))=\frac1n[t^{n-1}]h'(t)\ee^{nh(t)}
             =\frac1n[t^n]\ee^{nh(t)}.
\]
For the finite formula, rewrite \eqref{qb:eq:model} as
$v=\Li_s(r\ee^v)/(u-v/2)$, insert an auxiliary multiplier, and apply
Lagrange inversion in that multiplier. The coefficient of $r^n$ in the
$k$th power of the numerator is $C^{(s)}_{n,k}\ee^{nv}$ and vanishes for
$k>n$. Expanding $(u-v/2)^{-k}$ and the exponential gives
\eqref{qb:eq:finite-coeff}. Positivity also follows directly from either
of the last two expressions.
\end{proof}

\begin{corollary}[An exact positive coefficient identity]\label{qb:cor:prob}
If $s>2$ and $u\ge\Uc(s)$, then
\begin{equation}\label{qb:eq:probability}
 nR^na_n=[t^n]\exp\bigl(n(h(t)-h(1))\bigr).
\end{equation}
The right side is the probability that
\(\sum_{j\ge1}jN_j=n\), where the independent $N_j$ have Poisson means
$nh_j$. Their total mean count is $nh(1)<\infty$.
\end{corollary}
\begin{proof}
On $|t|\le1$, $2|\Li_s(t)|/u^2\le2A/u^2<1$. Thus the lift is well
defined and $h(1)=u-w=\vb$. Multiply \eqref{qb:eq:lagrange} by
$n\ee^{-n\vb}$. The exponential of
$n\sum h_j(t^j-1)$ is the stated probability generating function.
\end{proof}

This identity provides a positive numerical recurrence without subtracting
large alternating terms. It is also an exact bridge to classical random
allocation methods; no probability limit theorem is needed for the proofs
of the fixed-parameter laws below.

\section{Global control up to the boundary}\label{qb:sec:global}

\begin{proposition}[No hidden singularity inside the proposed circle]
\label{qb:prop:global}
Suppose $s>2$ and $u\ge\Uc(s)$. Then the germ $v$ extends analytically to
$|r|<R$, its Taylor series converges absolutely on $|r|\le R$, and
\begin{equation}\label{qb:eq:sum}
 \sum_{n\ge1}a_nR^n=\vb.
\end{equation}
It extends analytically through every point of $|r|=R$ except possibly
$r=R$.
\end{proposition}
\begin{proof}
For $0\le t\le1$, put
\[
 v(t)=u-\sqrt{u^2-2\Li_s(t)},\qquad r(t)=t\ee^{-v(t)}.
\]
The square root is at least $w\ge B>0$. On $0<t<1$,
\begin{equation}\label{qb:eq:monotone}
 \frac{\dd}{\dd t}\log r(t)=\frac1t
 \left(1-\frac{\Li_{s-1}(t)}{\sqrt{u^2-2\Li_s(t)}}\right)>0.
\end{equation}
Indeed $\Li_{s-1}(t)<B\le w$. Hence this parameterization supplies the
real solution on $0\le r<R$, tending to $\vb$ as $r\uparrow R$.

To pass from a real branch to its complex Taylor disk, iterate
\[
 v_0=0,\qquad
 v_{j+1}=\Phi(r,v_j),\qquad
 \Phi(r,v)=\frac{v^2}{2u}+\frac{\Li_s(r\ee^v)}u.
\]
All coefficients are nonnegative and increase coefficientwise. Each fixed
coefficient stabilizes: if two successive iterates agree modulo $r^k$,
their images agree modulo $r^{k+1}$, since $\Phi_v(0,0)=0$ and both
iterates vanish at zero. For each fixed $0\le r<R$, the iterates are
bounded by the real solution, since $\Phi$ preserves order on the interval
between zero and that solution. Monotone convergence of the coefficients
therefore proves convergence of the full formal series at every such $r$.
Its sum agrees with the real solution near zero, and then on the full
interval by analytic continuation. Taking $r\uparrow R$ proves
\eqref{qb:eq:sum}. Nonnegative coefficients imply absolute and uniform
convergence on the closed disk.

Let $r_0\ne R$ lie on $|r|=R$. Since every $a_n$ is positive, strict
triangle inequality gives $|v(r_0)|<\vb$. To justify strictness, equality
would force all $r_0^n/R^n$, $n\ge1$, to have the same argument, which the
terms $n=1,2$ already exclude unless $r_0=R$. Consequently
\[
 |r_0\ee^{v(r_0)}|=R\ee^{\Re v(r_0)}<R\ee^{\vb}=1.
\]
The equation is analytic there, and
\begin{align*}
 |\Phi_v(r_0,v(r_0))|
 &\le\frac{|v(r_0)|+\Li_{s-1}(|r_0\ee^{v(r_0)}|)}u\\
 &<\frac{\vb+B}{u}\le1.
\end{align*}
The analytic implicit function theorem gives continuation through $r_0$.
This proves the proposition without inferring a complex radius from the
real parameterization alone.
\end{proof}

The local singular expansions proved below are valid in a sector at $R$
with opening larger than $\pi$. Together with the proposition, compactness
of the remaining circle supplies the usual dented-disk, or
$\Delta$-domain, continuation. A nonzero local singular term then makes
$R$ the exact radius and the unique dominant singularity. This is the
analytic justification required before applying coefficient transfer.

\section{Logarithmic coordinates and a transfer dictionary}\label{qb:sec:coordinates}

Near the boundary use the determinations real on $0<r<R$:
\begin{equation}\label{qb:eq:coordinates}
 X=-\log(r/R),\qquad \tau=-\log t,\qquad
 v=\vb+X-\tau.
\end{equation}
The exact inverse relation is
\begin{equation}\label{qb:eq:H}
 X=H(\tau):=\tau-
 \left(\sqrt{u^2-2\Li_s(\ee^{-\tau})}-w\right).
\end{equation}
For $s\notin\N$, the classical convergent expansion is
\begin{equation}\label{qb:eq:jonquiere}
 \Li_s(\ee^{-\tau})=
 \Gamma(1-s)\tau^{s-1}+
 \sum_{k\ge0}\frac{(-1)^k\zeta(s-k)}{k!}\tau^k,
 \qquad |\tau|<2\pi,\quad |\arg\tau|<\pi.
\end{equation}
Its source is DLMF 25.12.12 \cite{dlmfpoly}. At a positive integer $m$,
the cancelling Gamma and zeta poles give
\begin{equation}\label{qb:eq:integer}
 \Li_m(\ee^{-\tau})=
 \frac{(-\tau)^{m-1}}{(m-1)!}\bigl(H_{m-1}-\log\tau\bigr)
 +\sum_{\substack{k\ge0\\k\ne m-1}}
 \frac{(-1)^k\zeta(m-k)}{k!}\tau^k,
\end{equation}
where $H_j=\sum_{k=1}^j1/k$. This follows either by taking the cancelling
limit in \eqref{qb:eq:jonquiere}, or by differentiating in $\tau$ and
matching the elementary $m=1$ identity.

Expanding the square in \eqref{qb:eq:model} after
\eqref{qb:eq:coordinates}, rather than prematurely expanding an inverse,
gives the exact local equation, for noninteger $s$,
\begin{align}\label{qb:eq:master}
 wX={}&(w-B)\tau+\Gamma(1-s)\tau^{s-1}
 +\frac{1+\zeta(s-2)}2\tau^2-\tau X+\frac{X^2}{2}\notag\\
 &+\sum_{k\ge3}\frac{(-1)^k\zeta(s-k)}{k!}\tau^k.
\end{align}
The additional $1$ in the quadratic coefficient is the contribution of
$-v^2/2$. It must not be dropped when importing a linear-core formula.

\subsection{Why powers of \texorpdfstring{$X$}{X} have no spurious descendants}

\begin{lemma}[Polylogarithmic transfer coordinates]\label{qb:lem:transfer}
Let a function analytic at $r=0$ have a unique dominant point $R$ and a
$\Delta$-continuation there. Write $X=-\log(r/R)$ locally. A singular term
$cX^\lambda$, $\lambda\notin\N_0$, contributes
\begin{equation}\label{qb:eq:power-transfer}
 \frac{c}{\Gamma(-\lambda)}R^{-n}n^{-\lambda-1}
\end{equation}
to coefficient asymptotics, without an additional string of integer
$1/n$ descendants from that term alone. A term $cX^k$, $k\in\N_0$, is
locally analytic and contributes no algebraic singularity term. A term
$cX^k\log X$ contributes
\begin{equation}\label{qb:eq:integer-transfer}
 c(-1)^{k+1}k!\,R^{-n}n^{-k-1}.
\end{equation}
More generally, for fixed $j$ the contribution of
$X^\lambda(\log X)^j$ is obtained by applying $\partial_\lambda^j$ to
$R^{-n}n^{-\lambda-1}/\Gamma(-\lambda)$.

These statements concern local contributions to an asymptotic expansion;
$X^\lambda$ need not itself be an analytic germ at $r=0$.
\end{lemma}
\begin{proof}
Put $z=r/R$. For $\lambda\notin\N_0$, the convergent identity
\eqref{qb:eq:jonquiere} gives
\[
 \Li_{\lambda+1}(z)=\Gamma(-\lambda)(-\log z)^\lambda+
 \text{a function analytic near }z=1.
\]
Subtract $c\Li_{\lambda+1}(z)/\Gamma(-\lambda)$ from the original
function. The subtracted function has the exact $n$th coefficient
\[c n^{-\lambda-1}/\Gamma(-\lambda).\]
The indicated singular term has therefore been removed. A finite collection of such subtractions reduces the
remainder to the ordinary big-$O$ transfer theorem \cite{fo,fs}; any local
analytic part can first be removed to the required order by a polynomial
in $1-z$. This proves the absence of artificial descendants.

For integer $k$, \eqref{qb:eq:integer} gives the same conclusion using
$\Li_{k+1}$. Equivalently,
\[
 \left.\frac{\dd}{\dd\lambda}\frac1{\Gamma(-\lambda)}
 \right|_{\lambda=k}=(-1)^{k+1}k!.
\]
For logarithmic powers, differentiate the normally convergent local
identity in a parameter neighborhood not containing a pole, or use its
pole-cancelled integer limit. Local analytic terms remain irrelevant to
any fixed algebraic asymptotic order. This proves the finite-order
statement, not convergence of an infinite coefficient-asymptotic series.
\end{proof}

\begin{remark}[Cancellation versus a different coordinate]
Expanding $X^\lambda$ first in powers of $1-r/R$ produces many terms.
Applying ordinary power transfer to those terms separately is legitimate
only after combining all terms of the desired order. The apparent integer
$1/n$ corrections cancel. Lemma~\ref{qb:lem:transfer} performs this
cancellation exactly by using a polylogarithm as the transfer model.
\end{remark}

\subsection{Local inverse sectors}

All subsequent local estimates are uniform on each closed subsector
$|\arg X|\le\pi-\epsilon$, $0<|X|<\delta_\epsilon$. For a simple leading
relation $X=c\tau^p(1+o(1))$, with $c>0$ and $p\ge1$, the selected inverse
is obtained near $\tau=(X/c)^{1/p}$, real positive for $X>0$.
One may prove sectorial inversion by Rouch\'e's theorem on a small relative
circle about this root: the leading map has one simple root there,
and the perturbation and its first derivative are smaller uniformly.
The root depends holomorphically on $X$ away from zero, and uniqueness
glues the local choices. The explicit implicit-function charts below give
a stronger convergent construction when $s$ is noninteger. The logarithmic
case is treated separately with its Lambert core.

\section{The strict boundary theorem, including integer orders}\label{qb:sec:strict}

\begin{theorem}[Proof of the pending boundary law]\label{qb:thm:strict}
Let $s>2$ and $u>\Uc(s)$. Set
\(\lambda=1-B/w>0\). Then the exact radius and unique dominant point are
$R=\ee^{-(u-w)}$, and
\begin{equation}\label{qb:eq:strict}
 a_n=\frac1{w\lambda^s}R^{-n}n^{-s}(1+o(1)).
\end{equation}
For noninteger $s$, the leading nonanalytic part of $v$ in $X$ coordinates is
\begin{equation}\label{qb:eq:strict-sing}
 \frac{\Gamma(1-s)}{w\lambda^s}X^{s-1}.
\end{equation}
For integer $m\ge3$, it is
\begin{equation}\label{qb:eq:strict-int}
 \frac{(-1)^m}{(m-1)!\,w\lambda^m}X^{m-1}\log X.
\end{equation}
In each case the omitted local terms are analytic or have strictly higher
singular order.
\end{theorem}
\begin{proof}
First suppose $s$ is noninteger. Remove the fractional term from
\eqref{qb:eq:jonquiere}, leaving an analytic function $Z(\tau)$, and let
\[
 H_0(\tau)=\tau-
 \left(\sqrt{u^2-2Z(\tau)}-w\right).
\]
Then $H_0(0)=0$, $H_0'(0)=\lambda$, and its ordinary analytic inverse
$\tau_0(X)$ satisfies $\tau_0(X)=X/\lambda+O(X^2)$.
Expanding the square root about its nonzero value $w$ gives
\[
 H(\tau)=H_0(\tau)+\frac{\Gamma(1-s)}w\tau^{s-1}
 +O(\tau^s)+O(\tau^{2s-2})
\]
on the chosen sectors. Sectorial inversion, or substitution followed by
Rouch\'e's theorem, therefore gives
\[
 \tau(X)=\tau_0(X)-
 \frac{\Gamma(1-s)}{w\lambda^s}X^{s-1}
 +O(X^{s-1+\eta})
\]
for any fixed $0<\eta<\min(1,s-2)$. This bound permits analytic terms to
be absorbed in $\tau_0$; the nonlinear interaction of two leading
fractional corrections first has order $2s-3>s-1$.
The identity $v=\vb+X-\tau$ proves \eqref{qb:eq:strict-sing}.

At an integer $m\ge3$, use \eqref{qb:eq:integer} instead. Its analytic
part again has a simple inverse, and its leading nonanalytic term is
$\ell_m\tau^{m-1}\log\tau$, where
$\ell_m=(-1)^m/(m-1)!$. The same calculation gives
\eqref{qb:eq:strict-int}; terms involving $\log\lambda$ are analytic at
this order. The remaining nonanalytic terms are of higher order, with at
most fixed logarithmic factors, so they are negligible in the leading
coefficient asymptotic.

The nonzero singular terms show that $r=R$ is singular. Proposition
\ref{qb:prop:global} excludes a smaller radius and any other singularity
on that circle. The sectorial inverse and the compact continuation away
from $R$ supply a $\Delta$-domain. Lemma~\ref{qb:lem:transfer} now gives
\eqref{qb:eq:strict} in both cases. Finally
$1/(w\lambda^s)=w^{s-1}/(w-B)^s$, which is precisely
\eqref{qb:eq:conjecture}.
\end{proof}

\begin{remark}[The singular limit in the theorem]
The constant in \eqref{qb:eq:strict} diverges as $w\downarrow B$.
Thus substituting $u=\Uc(s)$ in this formula is not valid. Criticality
changes the inverse equation's leading order, not merely its prefactor.
\end{remark}

\section{The fractional critical interval: a convergent Hahn chart}
\label{qb:sec:stable}

Assume $u=\Uc(s)$ and $2<s<3$. Write
\begin{equation}\label{qb:eq:stable-constants}
 \alpha=s-1\in(1,2),\qquad G=\Gamma(-\alpha)>0,\qquad
 d=(B/G)^{1/\alpha},\qquad T=dX^{1/\alpha}.
\end{equation}
Set $\rho=T^{\alpha-1}$ and $\sigma=T^{2-\alpha}$; thus
$T=\rho\sigma$. Both generators tend to zero.

\begin{theorem}[Convergent two-generator critical inverse]
\label{qb:thm:hahn}
There is a unique analytic function $U(\rho,\sigma)$ near $(0,0)$ with
$U(0,0)=1$ such that the boundary inverse is
\begin{equation}\label{qb:eq:hahn-inverse}
 \tau=T U(\rho,\sigma),\qquad v=\vb+X-TU(\rho,\sigma).
\end{equation}
It is determined by the analytic equation
\begin{align}\label{qb:eq:hahn-E}
 1={}&U^\alpha+\frac{1+\zeta(\alpha-1)}{2G}\sigma U^2
 -\frac{\rho\sigma}{B}U+\frac{G}{2B^2}\rho^2\sigma\notag\\
 &+\frac{\sigma}{G}\sum_{k\ge3}
 \frac{(-1)^k\zeta(\alpha+1-k)}{k!}
 (\rho\sigma)^{k-2}U^k.
\end{align}
Its Taylor series gives an absolutely convergent generalized power series
for $\tau$, with support contained in
\begin{equation}\label{qb:eq:stable-grid}
 \left\{\frac{1+j(\alpha-1)+k(2-\alpha)}{\alpha}:
                  j,k\in\N_0\right\}.
\end{equation}
Only finitely many support exponents lie below any fixed bound; collisions
are combined by addition.

If $\alpha=p/q$ in lowest terms, the exact minimum ramification index is
$p$: the inverse is a convergent Puiseux series in $X^{1/p}$, and cannot
be expressed as a convergent power series in $X^{1/M}$ with a smaller
positive integer $M$. If $\alpha$ is irrational, no finite Puiseux
ramification is possible.
\end{theorem}
\begin{proof}
At criticality $w=B$. In \eqref{qb:eq:master} put
$X=(G/B)T^\alpha$ and $\tau=TU$, and divide by $GT^\alpha$.
The identities
$T^{2-\alpha}=\sigma$, $T=\rho\sigma$, and
$T^\alpha=\rho^2\sigma$ give exactly \eqref{qb:eq:hahn-E}.
The tail is analytic near $(\rho,\sigma,U)=(0,0,1)$ because
\eqref{qb:eq:jonquiere} has a convergent analytic tail. Choose the power
$U^\alpha$ analytic near $U=1$. The derivative of the right side with
respect to $U$ is $\alpha\ne0$ at the base point. The analytic implicit
function theorem therefore produces $U$ and an absolutely convergent
bivariate Taylor series on a polydisk.

Substitution of the two small generators gives absolute convergence on
closed subsectors and the support \eqref{qb:eq:stable-grid}. Its
nonnegative integer generators have positive weights, proving local
finiteness. The branch agrees with the real inverse for $X>0$ small,
by its leading behavior, and hence with the continued germ selected in
Section~\ref{qb:sec:global}.

For $\alpha=p/q$, the exponents are
$\{(q+j(p-q)+k(2q-p))/p\}$. They are nonnegative integer multiples
of $1/p$, and absolute convergence gives an analytic series in $X^{1/p}$.
The leading term $dX^{q/p}$ is nonzero, so any possible ramification
index must be divisible by $p$. For irrational $\alpha$, its leading
exponent $1/\alpha$ is irrational and already excludes finite
ramification.
\end{proof}

\subsection{A finite formula for every chart coefficient}

The preceding implicit function is not merely existential. The following
finite recipe applies to both critical charts in this paper.

\begin{lemma}[Finite bivariate reversion]\label{qb:lem:finite-chart}
Suppose $E(p,q,U)=E_0(U)+\Delta E(p,q,U)$ is analytic near $(0,0,1)$,
$E_0(1)=0$, $E_0'(1)\ne0$, and $\Delta E(0,0,U)=0$.
Let $U(0,0)=1$ solve $E=0$, and define
\[
 Q(z)=\frac{E_0(1+z)}z,\qquad
 \Delta(p,q,z)=\Delta E(p,q,1+z).
\]
For $j+k\ge1$,
\begin{equation}\label{qb:eq:finite-chart}
 [p^jq^k](U-1)=
 \sum_{\ell=1}^{j+k}\frac{(-1)^\ell}{\ell}
 [p^jq^kz^{\ell-1}]
 \left(\frac{\Delta(p,q,z)}{Q(z)}\right)^\ell.
\end{equation}
Thus each coefficient requires only finite Taylor data.
\end{lemma}
\begin{proof}
Writing $z=U-1$ converts the equation to $z=-\Delta(p,q,z)/Q(z)$.
Insert a multiplier $a$ on the right and apply the ordinary Lagrange
formula in $a$. Its $\ell$th coefficient is
$(-1)^\ell[ z^{\ell-1}](\Delta/Q)^\ell/\ell$.
Since $\Delta$ has positive total $(p,q)$ degree, terms with
$\ell>j+k$ cannot contribute to $p^jq^k$. Setting $a=1$ is therefore
legitimate coefficientwise, giving the finite formula. Analyticity follows
independently from the implicit function theorem.
\end{proof}

For \eqref{qb:eq:hahn-E}, take $E_0(U)=U^\alpha-1$ and let $\Delta E$
be its displayed parameter-dependent part. This makes every Hahn
coefficient a finite expression in Gamma values, zeta values, and rational
functions of $\alpha$ and $B$.

\begin{corollary}[Coefficient expansion to every finite order]
\label{qb:cor:hahn-transfer}
Combine the collisions in Theorem~\ref{qb:thm:hahn} and write
$\tau=\sum_\lambda c_\lambda X^\lambda$. For any positive $L$ outside
the grid and outside $\N_0$,
\begin{equation}\label{qb:eq:hahn-transfer}
 a_nR^n=-\sum_{\lambda<L}
 \frac{c_\lambda}{\Gamma(-\lambda)}n^{-\lambda-1}
 +O(n^{-L-1}).
\end{equation}
The reciprocal Gamma factor is zero at nonnegative integer $\lambda$.
In particular \eqref{qb:eq:summary-stable} holds, and $R$ is the exact
radius and unique dominant point.
\end{corollary}
\begin{proof}
Absolute convergence of the two-variable chart and positivity of the two
weights imply that truncating below $L$ leaves $O(X^L)$ on each closed
subsector. One way to see the tail bound is to fix a smaller polydisk,
apply its geometric Taylor majorant, and factor $|X|^L$ from all omitted
monomials. Local finiteness leaves a positive gap before the first omitted
exponent. The leading term $dX^{1/\alpha}$ is nonanalytic and nonzero.
Proposition~\ref{qb:prop:global} and the chart supply exact dominance and
$\Delta$-continuation. Apply Lemma~\ref{qb:lem:transfer} to the finite
truncation and then big-$O$ transfer to its remainder. The analytic term
$\vb+X$ contributes nothing at these algebraic orders.
\end{proof}

The expansion in \eqref{qb:eq:hahn-inverse} is genuinely convergent.
Equation~\eqref{qb:eq:hahn-transfer}, in contrast, is asserted to every
\emph{fixed finite order}; no convergence claim is made for its infinite
coefficient-asymptotic sum.

\subsection{An exact cancellation at order \texorpdfstring{$s=5/2$}{5/2}}

\begin{corollary}[The missing $n^{-2}$ term]\label{qb:cor:fivehalves}
At $s=5/2$, $u=\Uc(5/2)$, put
\begin{equation}\label{qb:eq:fivehalf-constants}
 B=\zeta(3/2),\quad G=\Gamma(-3/2)=\frac{4\sqrt\pi}{3},\quad
 d=(B/G)^{2/3},\quad c=\frac{1+\zeta(1/2)}{2G},\quad
 K_* =\frac{7c^2}{9}+\frac{2}{3B}.
\end{equation}
Then
\begin{equation}\label{qb:eq:fivehalf-law}
 a_nR^n=
 -\frac{d}{\Gamma(-2/3)}n^{-5/3}
 -\frac{K_*d^2}{\Gamma(-4/3)}n^{-7/3}
 +O(n^{-8/3}).
\end{equation}
The coefficient of $n^{-2}$ is exactly zero. The second displayed
coefficient is strictly negative.
\end{corollary}
\begin{proof}
Here $\rho=\sigma=T^{1/2}$. Writing $z=T^{1/2}$ in
\eqref{qb:eq:hahn-E}, its first two coefficient equations are
\[
 1=U^{3/2}+czU^2-z^2U/B+O(z^3).
\]
Solving them gives
\begin{equation}\label{qb:eq:fivehalf-U}
 U=1-\frac{2c}{3}T^{1/2}
      +\left(\frac{7c^2}{9}+\frac2{3B}\right)T+O(T^{3/2}).
\end{equation}
Thus
\[
 v=\vb+X-T+\frac{2c}{3}T^{3/2}-K_*T^2+O(T^{5/2}).
\]
Since $T=dX^{2/3}$, the term $T^{3/2}=d^{3/2}X$ is analytic. It is
annihilated by $1/\Gamma(-1)=0$ in the transfer dictionary, rather than
producing a coefficient term of order $n^{-2}$. The other terms give
\eqref{qb:eq:fivehalf-law}. Finally $K_*>0$ and
$\Gamma(-4/3)>0$.
\end{proof}

This is an arithmetic cancellation in a specified quadratic model, not a
new general discovery that analytic local terms vanish under singularity
transfer. The related mechanism is already present in the repository's
linear-core critical Hahn theory \cite{repocritical}.

\section{The marginal order three and its complete logarithmic hierarchy}
\label{qb:sec:marginal}

Now let $s=3$, $u=\Uc(3)$, and $B=\zeta(2)$. Pole cancellation in
\eqref{qb:eq:integer} gives
\begin{equation}\label{qb:eq:Li3-local}
 \Li_3(\ee^{-\tau})=\zeta(3)-B\tau+
 \frac{\tau^2}{2}\left(\frac32-\log\tau\right)
 +\frac{\tau^3}{12}+O(\tau^4).
\end{equation}
The quadratic equation therefore has
\begin{equation}\label{qb:eq:marginal-H}
 X=\frac{\tau^2}{2B}\left(\log\frac1\tau+\frac52\right)
       +O\left(\tau^3\log\frac1\tau\right).
\end{equation}
The constant is $5/2$, not $3/2$. The extra $1$ is the quadratic-core
contribution in \eqref{qb:eq:master}.

\begin{theorem}[The exact Lambert core]\label{qb:thm:Lambert}
Define the large positive solution $Y$ by
\begin{equation}\label{qb:eq:Y}
 Y-\log Y=\log(1/X)+5-\log(4B).
\end{equation}
For positive small $X$ this is equivalently
\begin{equation}\label{qb:eq:Lambert}
 Y=-W_{-1}(-4B\ee^{-5}X),\qquad
 \tau_0(X)=2\sqrt{\frac{BX}{Y}}.
\end{equation}
Then the true inverse satisfies, on the selected complex sectors,
\begin{equation}\label{qb:eq:Lambert-error}
 \tau(X)=\tau_0(X)+O(\tau_0(X)^2).
\end{equation}
The large solution of \eqref{qb:eq:Y}, continued from $X>0$, defines the
complex branch in this statement; a numerical Lambert branch index is not
to be used without that continuation convention.
\end{theorem}
\begin{proof}
Let $H_0(\tau)=\tau^2(\log(1/\tau)+5/2)/(2B)$. Setting
$Y=2\log(1/\tau_0)+5$ converts $H_0(\tau_0)=X$ into
$Y\ee^{-Y}=4B\ee^{-5}X$, proving \eqref{qb:eq:Lambert}.
The large solution is holomorphic for $\Re\log(1/X)$ sufficiently large
and bounded imaginary part: the derivative of $Y-\log Y$ is
$1-1/Y$, uniformly separated from zero. Its real determination is the
usual $W_{-1}$ branch \cite{dlmfw}.

Uniformly on closed subsectors,
$|H_0'(\tau_0)|\asymp|\tau_0\log\tau_0|$ and
$H(\tau)-H_0(\tau)=O(\tau^3\log\tau)$ by
\eqref{qb:eq:marginal-H}. On a circle
$|\tau-\tau_0|=C|\tau_0|^2$, the linear part of
$H_0(\tau)-H_0(\tau_0)$ dominates this perturbation when $C$ is chosen
large and then $X$ small. Higher Taylor terms are smaller. Rouch\'e's
theorem gives one root there; its derivative is nonzero. The roots glue
holomorphically and match the physical real inverse. This proves
\eqref{qb:eq:Lambert-error} and provides the required sectorial
continuation.
\end{proof}

The error in \eqref{qb:eq:Lambert-error} is power-smaller than the leading
square root. It therefore lies beyond every fixed inverse-logarithmic
order in the coefficient expansion. This is stronger than retaining just
the leading regularly varying factor.

\subsection{All coefficients by finite Lagrange and Gamma formulas}

Let $I(a)=1/\Gamma(-a)$ and define
\begin{equation}\label{qb:eq:gj}
 g_j=\frac{I^{(j)}(1/2)}{I(1/2)},\qquad g_0=1.
\end{equation}
For example,
\begin{equation}\label{qb:eq:g12}
 g_1=p:=\psi(-1/2)=2-\gamma-2\log2,\qquad
 g_2=p^2-\psi^{(1)}(-1/2),\qquad
 \psi^{(1)}(-1/2)=4+\frac{\pi^2}{2}.
\end{equation}
These are exact Gamma and polygamma constants, not fitted parameters.
Define polynomials $Q_k$ and $P_m$ by
\begin{align}
 Q_0(A)&=1,\notag\\
 Q_k(A)&=-\frac1{2k}[z^{k-1}]
       (1+z)^{-3/2}(A+\log(1+z))^k,
       &&k\ge1,\label{qb:eq:Q}\\
 P_m(A)&=\sum_{j=0}^m\frac{(-1)^j g_j}{j!}
 \left\{\prod_{\ell=0}^{j-1}
 \left(\frac{\dd}{\dd A}-m+j-\ell-\frac12\right)\right\}
 Q_{m-j}(A).\label{qb:eq:P}
\end{align}
The product is the identity when $j=0$.

\begin{theorem}[Every inverse-logarithmic coefficient]\label{qb:thm:log-all}
Let $L=\log n$ and $A_n=\log L+5-\log(4B)$. For each fixed integer
$M\ge0$,
\begin{equation}\label{qb:eq:all-log}
 a_nR^n=\sqrt{\frac{B}{\pi}}\,n^{-3/2}L^{-1/2}
 \left\{\sum_{m=0}^M\frac{P_m(A_n)}{L^m}
 +O\left(\frac{(\log L)^{M+1}}{L^{M+1}}\right)\right\}.
\end{equation}
Each $P_m$ has degree $m$, with leading coefficient
$\binom{-1/2}{m}$. In particular, with
\begin{equation}\label{qb:eq:D}
 D_n=\log\log n+3+\gamma-\log B,
\end{equation}
\begin{align}\label{qb:eq:log-two}
 a_nR^n={}&\sqrt{\frac{B}{\pi}}\,n^{-3/2}L^{-1/2}
 \left[1-\frac{D_n}{2L}\right.\notag\\
 &\left.\quad+\frac{3D_n^2/8-D_n/2-\frac38(4+\pi^2/2)}{L^2}
 +O\left(\frac{(\log L)^3}{L^3}\right)\right].
\end{align}
The leading term is nonanalytic, so $R$ is the exact radius and unique
dominant point.
\end{theorem}
\begin{proof}
Put $\mathcal L=\log(1/X)$, $c=5-\log(4B)$,
$A=\log\mathcal L+c$, and write $Y=\mathcal L H$ in
\eqref{qb:eq:Y}. Then
\begin{equation}\label{qb:eq:H-Lambert}
 H=1+\frac{A+\log H}{\mathcal L}.
\end{equation}
For fixed $A$, ordinary Lagrange inversion for $H-1$ gives
\[
 H^{-1/2}=\sum_{k\ge0}Q_k(A)\mathcal L^{-k}
\]
with \eqref{qb:eq:Q}. Along $A=\log\mathcal L+c$, this expansion to any
fixed order has remainder
$O(\mathcal L^{-M-1}(1+|A|)^{M+1})$. To justify uniformity, treat
$p=1/\mathcal L$ and $z=A/\mathcal L$ as independent variables in
$H=1+z+p\log H$. The solution is analytic near $(p,z)=(0,0)$;
its convergent double Taylor majorant supplies the asserted remainder
when both variables tend to zero. The same argument controls any fixed
number of derivatives.

For completeness, the finite logarithmic transfer calculation can be
made directly in $X$ coordinates. In a Cauchy coefficient integral use
$r=R\ee^{-X}$ and then $X=z/n$. The local singular integral has the
Hankel kernel $\ee^z z^{1/2}\dd z$ with no further $1/n$ Jacobian
corrections. If $f(\mathcal L)$ denotes a slowly varying factor, its
finite Taylor expansion at $L$ is
\[
 f(L-\log z)=\sum_{j=0}^J
 \frac{(-1)^j}{j!}f^{(j)}(L)(\log z)^j+\text{remainder}.
\]
The Hankel integral of $\ee^z z^{1/2}(\log z)^j$ divided by $2\pi\ii$
is $I^{(j)}(1/2)$. This is the reciprocal-Gamma Hankel formula and its
fixed derivatives. Thus the normalized transfer operator is
\begin{equation}\label{qb:eq:log-operator}
 f(L)\longmapsto
 \sum_{j\ge0}\frac{(-1)^jg_j}{j!}f^{(j)}(L),
\end{equation}
used only to the finite order required.

Here $f(\mathcal L)=\sum_kQ_k(\log\mathcal L+c)
\mathcal L^{-k-1/2}$. Differentiating one term $j$ times gives
\[
 \mathcal L^{-k-j-1/2}
 \prod_{\ell=0}^{j-1}
 \left(\frac{\dd}{\dd A}-k-\ell-\frac12\right)Q_k(A).
\]
Collect $k+j=m$ in \eqref{qb:eq:log-operator}; this is precisely
\eqref{qb:eq:P}. Multiplication by the singular prefactor
$-2\sqrt B$ and by $I(1/2)=-1/(2\sqrt\pi)$ gives
\eqref{qb:eq:all-log}.

The finite Hankel expansion has the stated uniform remainder. One can
split the scaled contour into $|\log z|\le L/2$ and its complement.
On the first part the analytic majorants above and the Taylor remainder
bound give the next inverse-logarithmic order after integration. Near
$z=0$, the factor $z^{1/2}$ makes the excluded portion exponentially small
in $L$; at infinity the tilted Hankel rays give exponential decay in
$|z|$. The part of the original contour away from $R$ is exponentially
small in $n$, by the $\Delta$-continuation. Finally the error
$O(\tau_0^2)=O(X/\mathcal L)$ from
Theorem~\ref{qb:thm:Lambert} contributes a power-smaller coefficient
error; even the weaker bound $O(X^{1-\epsilon})$, with
$0<\epsilon<1/2$, suffices to put it beyond every displayed logarithmic
order. No infinite asymptotic expansion has been interchanged with an
integral.

For the first terms,
\[
 Q_1=-A/2,\quad Q_2=3A^2/8-A/2,
\]
\[
 P_1=-A/2+p/2,\quad
 P_2=3A^2/8-3Ap/4-A/2+p/2+3g_2/8.
\]
Writing $D=A-p$ gives
$P_1=-D/2$ and
$P_2=3D^2/8-D/2-3\psi^{(1)}(-1/2)/8$, proving
\eqref{qb:eq:log-two}. Formula~\eqref{qb:eq:Q} shows that the highest
$A$ coefficient is $\binom{-1/2}{m}$; all $j>0$ terms in
\eqref{qb:eq:P} have lower degree. Exact dominance follows as before
from the nonzero singular term and Proposition~\ref{qb:prop:global}.
\end{proof}

The two elementary checks $g_1=2-\gamma-2\log2$ and
$\psi^{(1)}(-1/2)=4+\pi^2/2$ follow from Gamma recurrence, duplication,
and their derivatives \cite{dlmfgamma}. They explain why Euler's constant
and polygamma data appear even though the original equation contains only
a polylogarithm and a quadratic polynomial.

\section{Finite curvature does not always mean a Puiseux germ}
\label{qb:sec:finite}

For $s>3$ at criticality put
\begin{equation}\label{qb:eq:curvature}
 C=1+\zeta(s-2)>0,\qquad b=\frac{C}{2B}>0.
\end{equation}
Equation~\eqref{qb:eq:H} has $H(\tau)=b\tau^2(1+o(1))$.
This gives the leading square root and
\eqref{qb:eq:summary-finite}. It does not imply that the whole inverse is
an analytic function of $\sqrt X$.

\subsection{Every noninteger order above three}

\begin{theorem}[A second convergent critical chart]
\label{qb:thm:upper-chart}
Let $s>3$ be noninteger and $u=\Uc(s)$. Put
\[
 G=\Gamma(1-s),\qquad q=(X/b)^{1/2},\qquad \sigma=q^{s-3}.
\]
Then $\tau=qU(q,\sigma)$ for a unique analytic $U$ near $(0,0)$ with
$U(0,0)=1$. Its exact equation is
\begin{align}\label{qb:eq:upper-chart}
 1={}&U^2+\frac{G}{Bb}\sigma U^{s-1}
 -\frac qB U+\frac b{2B}q^2\notag\\
 &+\frac1{Bb}\sum_{k\ge3}
 \frac{(-1)^k\zeta(s-k)}{k!}q^{k-2}U^k.
\end{align}
The support in $X$ is contained in
\begin{equation}\label{qb:eq:upper-grid}
 \left\{\frac{1+j+k(s-3)}2:j,k\in\N_0\right\}.
\end{equation}
Lemma~\ref{qb:lem:finite-chart}, now with $E_0(U)=U^2-1$, gives every
coefficient by a finite formula. The expansion converges, and its finite
coefficient transfer uses Lemma~\ref{qb:lem:transfer} exactly as in
Corollary~\ref{qb:cor:hahn-transfer}.

Relative to the analytic square-root core obtained by deleting the Gamma
term in \eqref{qb:eq:jonquiere}, the first additional local term in $v$ is
\begin{equation}\label{qb:eq:upper-first}
 \frac{G}{2B}\,b^{-s/2}X^{(s-2)/2}.
\end{equation}
If $s=a/d$ in lowest terms is noninteger, the minimum Puiseux ramification
index is $2d$. If $s$ is irrational, no finite ramification exists.
\end{theorem}
\begin{proof}
Substitute $X=bq^2$ and $\tau=qU$ into
\eqref{qb:eq:master} with $w=B$, and divide by $Bbq^2$. Since
$Bb=C/2$, this is \eqref{qb:eq:upper-chart}. The equation is analytic
in $q,\sigma,U$ near $(0,0,1)$, and its $U$ derivative at that point is
$2$. The analytic implicit function theorem gives the claimed convergent
chart. Its two positive exponent weights give local finiteness of the
support and finite transfer to any required asymptotic order.

The coefficient of $\sigma$ in $U$ at $q=0$ is $-G/(2Bb)$. Therefore
the first term involving the new generator in $v=\vb+X-qU$ is
$Gq^{s-2}/(2Bb)$, which is \eqref{qb:eq:upper-first}. The remaining
terms with at least one $\sigma$ have strictly higher exponent; the
$\sigma=0$ chart is precisely the analytic square-root reference core.

For rational $s=a/d$, the grid is contained in $(2d)^{-1}\N$ and gives
an analytic series in $X^{1/(2d)}$. The nonzero leading square root and
the nonzero exponent $(s-2)/2$ in \eqref{qb:eq:upper-first} force the
ramification index to be divisible by both $2$ and the denominator of
$(a-2d)/(2d)$. Their least common multiple is $2d$. The extra exponent
cannot cancel with a pure half-integer term because $s$ is noninteger.
For irrational $s$, that exponent is irrational and nonzero, proving
nonramifiability.
\end{proof}

The leading local term is $-b^{-1/2}X^{1/2}$. Its coefficient transfer is
\[
 \frac{b^{-1/2}}{2\sqrt\pi}R^{-n}n^{-3/2}
 =\sqrt{\frac{B}{2\pi(1+\zeta(s-2))}}R^{-n}n^{-3/2},
\]
which proves \eqref{qb:eq:summary-finite} at every noninteger $s>3$.
It also establishes exact dominance by the global proposition.

\subsection{Integer orders retain a logarithmic obstruction}

\begin{theorem}[First resonant logarithm above the marginal order]
\label{qb:thm:integer-obstruction}
Let $m\ge4$ be an integer and $u=\Uc(m)$. Write
\[
 B=\zeta(m-1),\quad b=\frac{1+\zeta(m-2)}{2B},\quad
 \ell_m=\frac{(-1)^m}{(m-1)!}.
\]
There is a function $V(q)$ analytic near $q=0$, with
$V(\sqrt X)=\vb-b^{-1/2}X^{1/2}+O(X)$, such that
\begin{align}\label{qb:eq:integer-obstruction}
 v={}&V(\sqrt X)+\frac{\ell_m}{4B}b^{-m/2}
       X^{(m-2)/2}\log X\notag\\
 &+O\left(X^{(m-1)/2}(1+|\log X|^2)\right).
\end{align}
The displayed logarithmic coefficient is nonzero. Consequently no finite
Puiseux ramification represents the full inverse, despite its leading
square root. The leading coefficient law is still
\eqref{qb:eq:summary-finite}.
\end{theorem}
\begin{proof}
Write \eqref{qb:eq:integer} as
$\Li_m(\ee^{-\tau})=Z_m(\tau)+\ell_m\tau^{m-1}\log\tau$, with
$Z_m$ analytic. The corresponding analytic part of $H$ has a double zero
with leading coefficient $b$, so its inverse $\tau_a(X)$ is a convergent
series in $\sqrt X$, beginning with $b^{-1/2}X^{1/2}$.
Expanding the nonzero square root in \eqref{qb:eq:H} gives
\[
 H(\tau)-H_a(\tau)=\frac{\ell_m}{B}\tau^{m-1}\log\tau
 +O\left(\tau^m(1+|\log\tau|^2)\right).
\]
Inverting around $\tau_a$, whose forward derivative is
$2b\tau_a(1+O(\tau_a))$, gives
\[
 \tau-\tau_a=-\frac{\ell_m}{2Bb}\tau_a^{m-2}\log\tau_a
 +O\left(\tau_a^{m-1}(1+|\log\tau_a|^2)\right).
\]
The possible second perturbation term has order
$\tau_a^{2m-5}\log^2\tau_a$, which is no larger than the stated
remainder for $m\ge4$. Now
$\log\tau_a=\tfrac12\log X+\log(\tau_a/\sqrt X)$; the latter
factor is analytic in $\sqrt X$. Absorb its leading analytic contribution
into $V$ and use $v=\vb+X-\tau$. This yields
\eqref{qb:eq:integer-obstruction}.

The coefficient of the first logarithm cannot be removed by a function
analytic in any finite root of $X$: on the common finite cover all such
functions have ordinary power series, whereas the first nonzero logarithm
still has nontrivial logarithmic continuation. Equivalently its asymptotic
expansion has a logarithmic monomial absent from every Puiseux series.
The leading square root, sectorial inversion, and the global proposition
prove the coefficient law and unique dominance.
\end{proof}

\begin{corollary}[A nonzero $n^{-2}$ correction at order four]
\label{qb:cor:four}
At $s=4$, $u=\Uc(4)$, with $B=\zeta(3)$ and $C=1+\zeta(2)$,
\begin{equation}\label{qb:eq:four-law}
 a_nR^n=\sqrt{\frac{B}{2\pi C}}\,n^{-3/2}
       +\frac{B}{6C^2}\,n^{-2}
       +O(n^{-5/2}(\log n)^2).
\end{equation}
\end{corollary}
\begin{proof}
The logarithmic coefficient in
\eqref{qb:eq:integer-obstruction} is
$\ell_4b^{-2}/(4B)=B/(6C^2)$. Its term is $X\log X$ and transfers
with multiplier $1$ by \eqref{qb:eq:integer-transfer}. The analytic
$X$ term does not contribute. The remaining Puiseux terms and logarithmic
remainder have the stated coefficient order.
\end{proof}

Together the stable, marginal, and finite-curvature arguments finish the
proof of Theorem~\ref{qb:thm:summary}. They also show why it is necessary
to distinguish a leading square-root law from an analytic fold: at the
polylogarithm boundary the forward function itself need not be analytic.

\section{Finite-prefix changes and a normalization dictionary}
\label{qb:sec:prefix}

The pure polylogarithm is not essential to the boundary proof. Finite
changes in the low-index coefficients can be retained exactly without
changing the tail class. This is useful when comparing separate packages,
where a polynomial prefix or a different amplitude is often present.

Let
\begin{equation}\label{qb:eq:prefix-model}
 f_s(t)=\kappa\Li_s(t)+P(t),\qquad \kappa>0,\quad P(0)=0,
\end{equation}
where $P$ is a real polynomial. Assume that the final Taylor coefficients
$f_j$ are nonnegative and $f_1>0$. The coefficients of $P$ need not all be
nonnegative. Replace \eqref{qb:eq:model} by
\(uv-v^2/2=f_s(r\ee^v)\). Put $\Dop=t\,\dd/\dd t$ and define
\begin{align}\label{qb:eq:prefix-constants}
 A&=\kappa\zeta(s)+P(1),&
 B&=\kappa\zeta(s-1)+\Dop P(1),\notag\\
 \Uc&=\sqrt{2A+B^2},&
 w&=\sqrt{u^2-2A},\quad \vb=u-w,\quad R=\ee^{-\vb}.
\end{align}

\begin{theorem}[Finite-prefix boundary transport]
\label{qb:thm:prefix}
For $s>2$, $u\ge\Uc$, the exact radius is $R$, the only dominant
singularity is $R$, and $\sum a_nR^n=\vb$. If $u>\Uc$, then
\begin{equation}\label{qb:eq:prefix-strict}
 a_nR^n\sim\frac{\kappa}{w(1-B/w)^s}n^{-s}.
\end{equation}
At $u=\Uc$, the fractional law for $2<s<3$ is
\begin{equation}\label{qb:eq:prefix-stable}
 a_nR^n\sim
 \frac{(B/(\kappa\Gamma(1-s)))^{1/(s-1)}}
      {-\Gamma(-1/(s-1))}n^{-1-1/(s-1)}.
\end{equation}
For $s>3$, put
$C=1+\kappa\zeta(s-2)+\Dop^2P(1)>0$. Then
\begin{equation}\label{qb:eq:prefix-finite}
 a_nR^n\sim\sqrt{\frac{B}{2\pi C}}n^{-3/2}.
\end{equation}
The convergent noninteger charts remain valid with
$G=\kappa\Gamma(1-s)$ and with each regular coefficient
$\zeta(s-k)$ replaced by
$\kappa\zeta(s-k)+\Dop^kP(1)$.

At the marginal order $s=3$, define
\begin{equation}\label{qb:eq:prefix-marginal}
 K=\frac B\kappa,\qquad
 b_0=\frac32+\frac{1+\Dop^2P(1)}\kappa.
\end{equation}
The exact Lambert core is
\begin{equation}\label{qb:eq:prefix-Lambert}
 Y=-W_{-1}(-4K\ee^{-2b_0}X),\qquad
 \tau_0=2\sqrt{\frac{KX}{Y}},
\end{equation}
with the same continuation convention as in
Theorem~\ref{qb:thm:Lambert}. The complete formula
\eqref{qb:eq:all-log} holds with prefactor $\sqrt{K/\pi}$ and
$A_n=\log\log n+2b_0-\log(4K)$. Its first two corrections are
\eqref{qb:eq:log-two} with
\begin{equation}\label{qb:eq:prefix-D}
 D_n=\log\log n+2b_0-2+\gamma-\log K.
\end{equation}
\end{theorem}
\begin{proof}
The lift is $h=u-\sqrt{u^2-2f_s(t)}$, with
$uh_j=f_j+\tfrac12\sum h_kh_{j-k}$. Nonnegative coefficients and
$f_1>0$ imply $h_j\ge0$, $h_1>0$, and $a_n>0$ for every $n$ by
\eqref{qb:eq:lagrange}; the term using $h_1^n$ already proves strictness.
Also $A>0$ and $B>0$. The proof of
Proposition~\ref{qb:prop:global} applies with $\Li_s$ replaced by
$f_s$ and $\Li_{s-1}$ by $\Dop f_s$. In particular
$\Dop f_s(t)<B$ for $0<t<1$, and the strict circle argument is unchanged.

The regular Taylor coefficients of $P(\ee^{-\tau})$ are
$(-1)^k\Dop^kP(1)/k!$, whereas the singular amplitude is multiplied by
$\kappa$. This proves all noninteger statements by the same exact
substitutions and implicit charts. At integer orders use the cancelled
formula \eqref{qb:eq:integer}. When $s>3$, $C=1+\sum j^2f_j>0$,
so the curvature is nonzero.

At $s=3$ the master equation has leading form
\[
 BX=\frac{\kappa\tau^2}{2}
 \left(\log\frac1\tau+\frac32+
                  \frac{1+\Dop^2P(1)}\kappa\right)
 +O(\tau^3\log(1/\tau)).
\]
This is $X=\tau^2(\log(1/\tau)+b_0)/(2K)$ plus the same
power-smaller remainder. Solving it gives
\eqref{qb:eq:prefix-Lambert}; the finite Lambert and Hankel calculations
are universal once $K$ and $b_0$ are substituted. This proves
\eqref{qb:eq:prefix-D} and the all-order claim.
\end{proof}

For the pure model, $\kappa=1$ and $P=0$, so $K=B$ and $b_0=5/2$.
For the numerical prefix test in the package,
$f_s(t)=2\Li_s(t)+t/2$, one has at $s=3$
\[
 B=2\zeta(2)+\tfrac12,\qquad K=\zeta(2)+\tfrac14,
 \qquad b_0=\tfrac94.
\]
Thus neither the radius nor the logarithmic shift can be copied from the
pure model without adjustment, even though the leading logarithmic
exponent is unchanged.

\section{Verification and reproducible evidence}\label{qb:sec:verification}

The program \path{code/verify.py} separates exact finite tests from
floating-point diagnostics. It writes only to an absent or empty output
directory, so a replay does not overwrite the delivered evidence.

\subsection{Exact checks}

The recorded run passed \textbf{293 exact assertions}. For
$s=2,3,4$, $u=1,2,3$, and $1\le n\le10$, rational arithmetic compares
three independent coefficient constructions: iteration of the original
quadratic fixed-point equation, the lift/Lagrange identity, and the
predecessor's positive finite formula. Positivity is checked at each
instance. Symbolic tests additionally compare the Lambert polynomials
$Q_0,\ldots,Q_4$ with a separate implicit recurrence, and the Hankel
polynomials $P_0,\ldots,P_4$ with direct substitution
$L\mapsto L-y$ followed by moment replacement. The first two shifted
corrections, the $s=5/2$ local coefficient equations, the $s=4$ logarithmic
balance, and the leading polynomial coefficients are also checked.

These finite computations are regressions, not proofs of the universal
claims. In particular, they do not replace the positive-circle or
sectorial-continuation arguments.

\subsection{Positive numerical coefficient recurrence}

For each target index $n$ define
\[
 p_k=[t^k]\exp(n(h(t)-h(1))).
\]
Then
\begin{equation}\label{qb:eq:prob-rec}
 p_0=\ee^{-n\vb},\qquad
 p_k=\frac nk\sum_{j=1}^k jh_jp_{k-j},\qquad
 a_nR^n=\frac{p_n}{n}.
\end{equation}
The stored diagnostics use this positive recurrence, not the asymptotic
formula, to compute the reference coefficients. Seven models were tested
at seven indices up to $4096$, for 49 coefficient evaluations. The
platform's NumPy \texttt{longdouble} has 18 decimal digits of precision;
these larger-index computations were \emph{not} performed at 85 digits.
For each model an independent arbitrary-precision recurrence at $n=32$
agreed with the extended-precision result to relative error below
$1.3\cdot10^{-18}$. The arbitrary-precision constants and local coordinate
checks used \texttt{mpmath} at 85 decimal digits.

On platforms whose \texttt{longdouble} exponent range is too small for
$p_0$, the code automatically uses an arbitrary-exponent \texttt{mpmath}
recurrence. This fallback is slower but avoids silently underflowing the
initial probability. Neither implementation supplies interval enclosures.

\begin{table}[htbp]\centering
\caption{Coefficient divided by its proved leading asymptotic. Strict
boundary tests use $w=2\zeta(s-1)$, hence $1-B/w=1/2$.}
\label{qb:tab:leading}
\begin{tabular}{lrrrr}\toprule
Regime & $n=128$ & $n=512$ & $n=2048$ & $n=4096$\\\midrule
Critical $s=5/2$ & 0.986604 & 0.994512 & 0.997779 & 0.998591\\
Critical $s=3$ & 0.655458 & 0.703836 & 0.739120 & 0.753441\\
Critical $s=4$ & 1.006135 & 1.003630 & 1.002010 & 1.001472\\
Strict, $s=5/2$ & 0.906424 & 0.973409 & 0.993100 & 0.996528\\
Strict, $s=3$ & 1.319718 & 1.092013 & 1.026750 & 1.014462\\
Strict, $s=4$ & 1.425213 & 1.071592 & 1.016618 & 1.008208\\\bottomrule
\end{tabular}
\end{table}

The marginal leading law approaches its limit slowly. Its corrections are
therefore useful rather than decorative. Table~\ref{qb:tab:refinement}
shows the coefficient divided by the refined approximation. In its
$s=3$ column both inverse-logarithmic corrections in
\eqref{qb:eq:log-two} are retained; the fractional column retains the
second nonzero term in \eqref{qb:eq:fivehalf-law}. The reference
coefficients in both columns still come from \eqref{qb:eq:prob-rec}.

\begin{table}[htbp]\centering
\caption{Independent coefficient divided by the refined approximation.}
\label{qb:tab:refinement}
\begin{tabular}{rrr}\toprule
$n$ & Critical $s=5/2$, two nonzero terms & Critical $s=3$, two corrections\\\midrule
128 & 1.001199056 & 1.049782812\\
512 & 1.000298579 & 1.020656682\\
2048 & 1.000074533 & 1.010075657\\
4096 & 1.000037253 & 1.007282765\\\bottomrule
\end{tabular}
\end{table}

There are 21 real or complex checks of the exact coordinate identity
against the original equation, with recorded absolute residuals below
$10^{-84}$, and six checks of the defining Lambert-core identity. Local
ratios at $\tau=10^{-4},10^{-8},10^{-12}$ test the predicted leading
normal forms; the data retain the observed errors rather than relabeling
them as certified bounds. The prefix example
$2\Li_3(t)+t/2$ tests the change of marginal normalization separately.
Full values and software versions are in \path{data/verification.json}
and the CSV files.

\section{Source audit, corrections, and integration}\label{qb:sec:audit}

\subsection{What was actually inspected}

The three user-specified areas were inspected through their GitHub
indexes and selected source material \cite{repopoly,repotrans,repoincoming}. The observed main commit was
\begin{center}\small
\texttt{6207c3d40047fd8d1400a4d509fee71c9a3a8132}.
\end{center}
Initial index reads were on \texttt{main}; the later linear-core TeX read
was pinned to this commit. Thus this is not falsely described as an atomic
full-repository checkout. The incoming directory and commit metadata list
five new polylogarithm/transseries archives. Their complete contents were
not all audited.

The predecessor \cite{prior} was read in a Library TeX copy named
\path{article(20261012-044824).tex}, specifically source lines 1--1000
and 1030--1487. Its exact title, equation, interior theorem, and pending
boundary question identify the research target. The Library copy was not
byte-compared with a member of
\path{docs/incoming/Resonant_Polylogarithm_Transseries.zip}; that
provenance limitation is explicit. The existing critical-Hahn package's
README and its TeX opening through source line 150 were also inspected
\cite{repocritical}. Its extensive editorial cross-references demonstrate
that generic critical/Hahn and marginal mechanisms already belong to the
repository's research program.

The external source check covered the DLMF polylogarithm formula and
Lambert material, the Janson publication record/abstract, and targeted
literature searches. Direct retrieval of the full Flajolet--Odlyzko
article was unsuccessful; it is cited as the classical transfer source,
not represented as newly audited line by line. The needed local transfer
and finite Hankel calculations are explained in this manuscript. Neither
the repository inspection nor the literature search establishes worldwide
publication priority.

\subsection{Status changes and safeguards, not invented source errors}

No new mathematical error was established in the accepted repository
statements inspected. In particular, the predecessor correctly labels
\eqref{qb:eq:conjecture} as unproved and correctly restricts its existing
fold theorem to the interior regime. The appropriate proposed editorial
change is to mark its boundary question as answered \emph{for fixed
parameters} by this article, subject to independent proof review. Its
joint-limit questions remain open.

There are four useful integration safeguards.
First, the simple-boundary formula cannot be continued by substitution to
criticality: its constant diverges. Second, inserting $t=1$ into the
interior curvature prefactor at $s=3$ would incorrectly erase the leading
coefficient; the missing factor is $(\log n)^{-1/2}$. Third, above $s=3$
a leading square root is not a theorem of analytic Puiseux structure:
Theorem~\ref{qb:thm:integer-obstruction} supplies an explicit nonzero
logarithmic obstruction at every integer $m\ge4$. Fourth, neglecting
Gamma zeros creates a false $n^{-2}$ contribution at $s=5/2$, whereas at
$s=4$ the logarithmic term really does produce such a contribution.
These are precise consequences of the present proofs, not accusations
about unseen source passages.

\subsection{Suggested placement and statement-level review}

A self-contained placement is
\begin{center}\small
\path{Analysis/Transseries/docs/series-and-transseries/}\par
\path{Quadratic_Polylogarithm_Boundary_Transfer/}.
\end{center}
The polylogarithm project index can link to this package, and the pending
article can receive a short editorial cross-reference rather than having
its original proofs silently rewritten. \path{INTEGRATION.md} and
\path{integration_note.tex} supply proposed wording. No remote change was
made.

A sensible review order is the exact coefficient identity and the master
equation, then the positive-circle argument, then the strict-boundary
proof and integer cancellation. The two convergent critical charts and
the Lambert/Hankel compiler can be reviewed independently afterward.
The proposed results do not establish general resurgence, Borel
summability, a uniform moving-critical-point law, or any new Lean theorem.

\section{Further research questions}\label{qb:sec:future}

\paragraph{1. Uniform approach to the critical surface.}
Theorem~\ref{qb:thm:strict} and the critical laws are fixed-parameter
statements. For $2<s<3$, balancing the drift term
$(w-B)\tau$ against $\Gamma(1-s)\tau^{s-1}$ in the master equation
suggests a coefficient crossover when
$w-B$ is of order $n^{-(s-2)/(s-1)}$. Prove a uniform coefficient profile,
including its normalizing constant and both outer limits. The balance is
a scaling suggestion, not a proved uniform theorem.

\paragraph{2. The joint order-three limit.}
Construct a pole-cancelled inverse chart for $s\to3$ and $n\to\infty$
that matches the two Hahn charts with the Lambert chart. The Gamma pole
and the pole of $\zeta(s-2)$ must remain paired. The linear-core marginal
and confluent packages are comparison sources; the quadratic drift terms
and normalization $b_0$ must be tracked explicitly.

\paragraph{3. The lower endpoint and large-core limits.}
As $s\downarrow2$, the moment $B=\zeta(s-1)$ diverges and the critical
surface escapes to infinity. Determine which joint scalings connect this
boundary theory to the dilogarithm's interior fold and to the predecessor's
large-core, large-sector question. Neither the fixed-$s$ formulas nor a
formal limit of their constants provides that connection.

\paragraph{4. Parameter derivatives and resonance cancellation.}
Differentiate the exact coefficient formulas with respect to $s$ and
$u$. Can one construct a uniform order-jet calculus for coefficients near
rational Hahn collisions and integer logarithmic resonance? The transfer
dictionary suggests explicit polygamma and $\log n$ terms, but uniform
control must precede differentiating an asymptotic error term.

\paragraph{5. Other exact cores.}
Replace $uv-v^2/2$ by the higher-degree exact polynomial cores of
integer-order Fermi inversion. The quadratic lift no longer applies
directly, and positivity may fail. Identify conditions on a general
analytic core that preserve the unique positive dominant boundary, and
determine how its first nonzero derivatives enter the marginal shift.

\paragraph{6. Periodic or signed tails.}
Replace the positive tail by a periodic-weight or cyclotomic
polylogarithm combination. Several boundary points can compete and the
strict triangle-inequality proof may break down. A finite Fourier
criterion for dominant-point multiplicity and cancellation would connect
this theory to the predecessor's periodic Borel-residue dictionary.

\paragraph{7. Effective convergence domains and error certificates.}
The analytic implicit-function proofs give genuine local convergence but
not optimal explicit polydisks. Produce interval-certified domains for
the two-variable charts and for the Lambert perturbation, and certified
coefficient enclosures for the positive recurrence. The floating-point
diagnostics delivered here are not such certificates.

\paragraph{8. Formalization and exact coefficient infrastructure.}
The finite reversion formula and the rational coefficient recurrence are
natural first proof-assistant targets. A second layer would formalize the
polylogarithmic transfer dictionary and its integer limit. The full radius
theorem additionally requires positive implicit-series convergence,
complex continuation, and Hankel transfer. The theorem manifest separates
these dependencies so that finite algebraic checks are not mistaken for
formal analytic coverage.

\section*{Conclusion}
\addcontentsline{toc}{section}{Conclusion}
The pending quadratic boundary law is an analytic theorem once positivity,
sectorial inversion, and integer-order pole cancellation are retained
together. The critical boundary carries more information than a leading
exponent: it has convergent Hahn charts away from integer resonance,
Gamma-controlled missing coefficient terms, and an exactly normalized
Lambert core at the marginal order. A finite quadratic term changes the
logarithmic constants and can produce resonant corrections invisible in a
naive square-root description. These refinements connect the exact
polylogarithm calculus to the repository's transseries inversion program
without claiming its classical or already-developed mechanisms as new.

\clearpage
\appendix
\section{Reproduction and artifact map}\label{qb:app:reproduce}

The article is self-contained in \path{article.tex}; compilation does not
require external images or font files. A TeX Live installation with the
standard packages named in the preamble is sufficient. The retained
numerical tables are rounded from the accompanying CSVs.
\begin{verbatim}
python -m pip install -r requirements.txt
python code/verify.py --output-dir replay --max-n 4096
pdflatex -interaction=nonstopmode -halt-on-error article.tex
pdflatex -interaction=nonstopmode -halt-on-error article.tex
pdflatex -interaction=nonstopmode -halt-on-error article.tex
\end{verbatim}
The output directory must be new or empty. The distributed
\path{data/verification.json} describes the executed run, not an
unperformed command. \path{BUILD_VALIDATION.json} records the PDF build
and page-layout checks. \path{THEOREMS.json} lists theorem dependencies
and the distinction between conventional proofs and finite checks.

\begin{longtable}{@{}>{\raggedright\arraybackslash}p{.31\textwidth}>{\raggedright\arraybackslash}p{.61\textwidth}@{}}
\toprule
Artifact & Purpose\\\midrule\endhead
\path{article.tex}, \path{article.pdf} & Complete proofs, exact identities,
research questions, and source distinctions.\\[4pt]
\path{code/verify.py} & Exact finite comparisons and positive numerical
coefficient diagnostics.\\[4pt]
\path{data/exact_coefficients.csv} & Rational coefficients for the finite
replay cases.\\[4pt]
\path{data/marginal_polynomials.txt} & $Q_0,\ldots,Q_4$ and
$P_0,\ldots,P_4$ in symbolic Gamma-moment coordinates.\\[4pt]
\path{data/coefficient_diagnostics.csv} & 49 positive-recurrence comparisons
through $n=4096$, with arithmetic engine recorded.\\[4pt]
\path{data/local_checks.csv}, \path{data/lambert_checks.csv} & Local
normal-form and defining-equation diagnostics.\\[4pt]
\path{SOURCES.md}, \path{INTEGRATION.md} & Provenance, scope limits, and
proposed filing/editorial notes.\\\bottomrule
\end{longtable}

\clearpage
\phantomsection
\addcontentsline{toc}{section}{References}
\begingroup\small
\begin{thebibliography}{12}
\setlength{\itemsep}{3pt}
\bibitem{prior}
Research draft prepared for Vladimir Reshetnikov with OpenAI assistance,
\emph{Resonant Polylogarithm Transseries: Gamma Borel Kernels, Cyclotomic
Order Jets, and Sharp Dilogarithm Reversion}, 11 October 2026.
Theorem 9.1 and Section 12, Question 1. Inspected Library TeX copy
\path{article(20261012-044824).tex}; associated incoming archive name
\path{docs/incoming/Resonant_Polylogarithm_Transseries.zip}.
Source ranges and provenance limitation are stated in Section~\ref{qb:sec:audit}.

\bibitem{dlmfpoly}
NIST Digital Library of Mathematical Functions,
\emph{\S25.12, Polylogarithms}, especially 25.12.10--25.12.12.
\url{https://dlmf.nist.gov/25.12}. Inspected 12 October 2026.
The Gamma--zeta expansion is attributed there to Erd\'elyi et al.,
\emph{Higher Transcendental Functions}, vol.~1 (1953), Section 1.11.

\bibitem{dlmfw}
NIST Digital Library of Mathematical Functions,
\emph{\S4.13, Lambert $W$-Function}, branch definitions and asymptotic
expansions. \url{https://dlmf.nist.gov/4.13}. Inspected 12 October 2026.

\bibitem{dlmfgamma}
NIST Digital Library of Mathematical Functions,
\emph{Chapter 5, Gamma Function}, especially recurrence, duplication,
and polygamma relations in \S\S5.5 and 5.15.
\url{https://dlmf.nist.gov/5}. Classical identities used in
\eqref{qb:eq:g12}; the finite recurrence evaluations are reproduced in
\path{code/verify.py}.

\bibitem{fo}
P. Flajolet and A. Odlyzko,
Singularity analysis of generating functions,
\emph{SIAM Journal on Discrete Mathematics} \textbf{3} (1990), 216--240.
\url{https://doi.org/10.1137/0403019}.
Classical transfer reference; direct full-text retrieval was unsuccessful
in the present audit.

\bibitem{fs}
P. Flajolet and R. Sedgewick,
\emph{Analytic Combinatorics}, Cambridge University Press, 2009,
Chapters VI and VII. Standard reference for singularity transfer and
positive implicit-function schemas; not represented as a new full-book
audit.

\bibitem{janson}
S. Janson,
Simply generated trees, conditioned Galton--Watson trees, random allocations
and condensation, \emph{Probability Surveys} \textbf{9} (2012), 103--252.
\url{https://arxiv.org/abs/1112.0510}.
Publication record and abstract inspected; used to identify the established
probabilistic framework, not as a substitute for the proofs above.

\bibitem{repocritical}
ProveIt research package,
\emph{Beyond Finite-Action Folds: Critical Hahn Transseries, Stable Sector
Asymptotics, and Sharp Action Budgets}, 29 September 2026, with editorial
amendments.
\path{Analysis/Transseries/docs/series-and-transseries/}
\path{Critical_Hahn_Transseries_Beyond_Finite_Action_Folds/}.
README and TeX source lines 1--150 inspected; the latter at commit
\texttt{6207c3d40047fd8d1400a4d509fee71c9a3a8132}.

\bibitem{repotrans}
ProveIt Contributors,
\emph{Series and transseries: canonical-volume and research-package index}.
\path{Analysis/Transseries/docs/series-and-transseries/README.md},
lines 1--220 inspected on 12 October 2026.

\bibitem{repopoly}
ProveIt Contributors,
\emph{Polylogarithms project index}.
\path{Analysis/Polylogarithms/README.md}, lines 1--190 inspected on
12 October 2026. Accepted/pending status is governed by the project ledger,
not by a historical article abstract alone.

\bibitem{repoincoming}
ProveIt Contributors,
\emph{Incoming research intake instructions and archive index}.
\path{docs/incoming/README.md}, lines 1--145, and directory/commit metadata
inspected on 12 October 2026.
\url{https://github.com/VladimirReshetnikov/ProveIt}.
\end{thebibliography}
\endgroup
\end{document}
